\section{Gauge fixing and the free-boundary state}\label{sec:state}

For a function $O$ of the links let $O^F$ be $O$ with every link of $F$ set to
the identity.

\begin{lemma}[gauge fixing]\label{lem:gaugefix}
Let $\mu^F_\Lambda$ be the probability measure on $G^{A\cap E(\Lambda)}$ with
density proportional to $\exp(N\beta\sum_{p\in P(\Lambda)}\ReTr Q^F_p)$. For
every gauge-invariant $O$ on $G^{E(\Lambda)}$, $W_\Lambda(O)=\mu^F_\Lambda(O^F)$.
\end{lemma}

\begin{proof}
Each segment of $F\cap E(\Lambda)$ is a path $x_0,\dots,x_k$. Given the
variables on $F$, put $g_{x_0}=1$ and $g_{x_{i+1}}=g_{x_i}Q_{(x_i,x_{i+1})}$
along every segment. The transformed field equals the identity on $F\cap
E(\Lambda)$, and for fixed values on $F$ the map from the active variables to
the transformed ones is a product of translations, which preserve Haar
measure. Gauge invariance of the action and of $O$ gives the claim.
\end{proof}

No holonomy remains because $F\cap E(\Lambda)$ is a forest. The measures
$\mu^F_\Lambda$ are finite-volume versions of the specification $\gamma^F$ on
$G^A$ in which $\gamma^F_e$ resamples an active $Q_e$ from the law with density
proportional to $\exp(N\beta\sum_{p\ni e}\ReTr Q^F_p)$ given the other active
links. Fix a bi-invariant metric $\delta$ on $G$ of diameter $D$ and define the
Kantorovich coefficients
\[
C_{ee'}=\sup\frac{W_1\bigl(\gamma^F_e(\cdot\mid Q),\gamma^F_e(\cdot\mid Q')\bigr)}
{\delta(Q_{e'},Q'_{e'})},
\]
the supremum over configurations that differ only at $e'$; $W_1$ is the
Kantorovich distance~\cite{Vil09}, and $C_{ee'}=0$ unless $e$ and $e'$ share a
plaquette.

\begin{lemma}[comparison]\label{lem:comparison}
Suppose every row and column of $(C_{ee'})$ sums to at most $\alpha<1$, and let
$\lambda>0$ with $\theta=\alpha e^\lambda<1$. Let $S\subset A$ and let $\mu_1,\mu_2$
be probability measures on $G^A$ invariant under $\gamma^F_e$ for $e\notin S$.
Let $T\subset A\setminus S$ be finite, $w_b=e^{-\lambda\dist(b,T)}$ with the
distance of link midpoints, and suppose $\sum_{z\in S}w_z<\infty$. There is a
coupling of $\mu_1$ and $\mu_2$ with $\varrho(b)=\E\,\delta(\sigma_b,\sigma'_b)$
satisfying
\[
\sum_{b\notin S}w_b\varrho(b)\le\frac{\theta D}{1-\theta}\sum_{z\in S}w_z ,
\]
and $\mu_1=\mu_2$ if $S=\emptyset$.
\end{lemma}

\begin{proof}
The couplings form a weakly compact set on which $\Psi(\pi)=\sum_bw_b\E_\pi
\delta(\sigma_b,\sigma'_b)$ is continuous; let $\pi^*$ minimise it. For $e\notin S$
replace $(\sigma_e,\sigma'_e)$ by a sample of an optimal coupling of
$\gamma^F_e(\cdot\mid\sigma)$ and $\gamma^F_e(\cdot\mid\sigma')$, chosen
measurably~\cite[Cor.~5.22]{Vil09}. Both $\mu_i$ are invariant, so this is again
a coupling, and changing the links that enter $\gamma^F_e$ one at a time gives
$\varrho(e)\le\sum_{e'}C_{ee'}\varrho(e')$ by minimality. Multiply by $w_e$ and sum
over $e\notin S$; where $C_{ee'}\ne0$, $w_e\le e^\lambda w_{e'}$, so
$\sum_{e\notin S}w_e\varrho(e)\le\theta[\sum_{e'\notin S}w_{e'}\varrho(e')+D\sum_{z
\in S}w_z]$. If $S=\emptyset$ this forces $\varrho\equiv0$.
\end{proof}

Conditioning on the links of the support of $H$~\cite[Ch.~3]{Kal21} and
applying the lemma to the two conditional measures gives, for the Gibbs
measure $\mu$ of $\gamma^F$, a local $F'$ that is $L$-Lipschitz in each link and
a bounded $H$ with finite support disjoint from that of $F'$,
\begin{equation}\label{eq:decay}
|\Cov_\mu(F',H)|\le\|H\|_\infty L\,\frac{\theta D}{1-\theta}\,|\operatorname{supp}H|\,
e^{-\lambda\dist(\operatorname{supp}F',\operatorname{supp}H)} .
\end{equation}

\begin{proposition}\label{prop:state}
Suppose every row and column of $(C_{ee'})$ sums to at most $\alpha<1$. Then
$\gamma^F$ has a unique Gibbs measure $\mu_F$; for every gauge-invariant local $O$
the limit $\omega(O)=\lim W_\Lambda(O)$ exists along every sequence of boxes
exhausting $\Z^4$ and equals $\mu_F(O^F)$; $\omega$ is invariant under
translations and the hypercubic group; truncated correlations of Lipschitz
gauge-invariant observables decay at every rate below $-\log\alpha$; and for
$\beta\ge0$, $\omega$ is reflection positive and its transfer matrix has
spectral gap at least $-\log\alpha$.
\end{proposition}

\begin{proof}
Gibbs measures exist by compactness and coincide by Lemma~\ref{lem:comparison}
with $S=\emptyset$. Extend $\mu^F_\Lambda$ by Haar measure outside $\Lambda$ and
let $S$ be the set of active links with a plaquette outside $P(\Lambda)$; off $S$
the extended measure is invariant under $\gamma^F_e$. Lemma~\ref{lem:comparison}
with $T$ the support of $O^F$ bounds $|\mu^F_\Lambda(O^F)-\mu_F(O^F)|$ by
$C_TR^3e^{-\lambda R}$, $R$ the distance from $T$ to the complement of $\Lambda$,
for Lipschitz $O$, and uniform approximation handles continuous $O$;
Lemma~\ref{lem:gaugefix} identifies $\mu^F_\Lambda(O^F)$ with $W_\Lambda(O)$. Since
$W_\Lambda(\tau_vO)=W_{\Lambda-v}(O)$ and the limit does not depend on the
exhausting sequence, $\omega$ is invariant; the line field is a device for
computing $W_\Lambda$ and is never moved. Reflection positivity of $W_\Lambda$ on
symmetric boxes passes to the limit. The decay is~\eqref{eq:decay}, since
$\Cov_\omega(O,O')=\Cov_{\mu_F}(O^F,O'^F)$. Lipschitz gauge-invariant functions
are dense in the gauge-invariant part of $L^2(\omega)$ on every support, being
Lipschitz functions on the compact quotient by the gauge group, so
Proposition~\ref{prop:gapdecay} gives the gap.
\end{proof}

\begin{proof}[Proof of Theorem~\ref{thm:A} for $\SU(2)$]
Identify $\SU(2)$ with the unit quaternions and use the chordal metric. For an
active link $e$ and a plaquette $p\ni e$ let $s_p$ be the product of the other
three variables, ordered so that $\operatorname{Re}(Q_es_p)=\frac12\Tr Q_p$, with
links of $F$ equal to one. Since $N\beta\ReTr Q_p=4\beta\ip{Q_e}{\bar s_p}$ in
$\R^4$, the conditional law of $Q_e$ is the von Mises--Fisher law $\nu_b(dx)
\propto e^{\ip xb}dx$ on $S^3$ with $b=4\beta\sum_{p\ni e}\bar s_p$. Moving a
partner $e'$ by $\delta$ moves one staple by an isometric image of that
displacement, hence $b$ by $4|\beta|\delta$, and the sharp
quarter (Theorem~\ref{q-thm:quarter} in Appendix~\ref{app:sharp}), $W_1(\nu_b,\nu_{b'})\le\frac14|b-b'|$, gives
$C_{ee'}\le|\beta|$. By Lemma~\ref{lem:structure}(iii) every row and column has
twelve nonzero entries, so $\alpha\le12|\beta|=3|\beta_W|$, and
Proposition~\ref{prop:state} applies for $|\beta_W|<\frac13$.
\end{proof}

The value $\frac13$ is the end of this argument: at $b=0$ the derivative of
$\nu_b$ has Kantorovich norm exactly $\frac14$, the six staples of an active
link can be arranged with sum zero while one of them moves freely, and so
every row sum of the matrix is exactly $3|\beta_W|$. Without the gauge every
link has eighteen partners and the same argument stops at $\beta_W=\frac29$
\cite[Prop.~13]{BB26a}.

